\documentclass[11pt,a4paper]{article}

\usepackage[margin=25mm]{geometry}
\usepackage{amsmath,amssymb,bm}
\usepackage{microtype}
\usepackage[hidelinks]{hyperref}

\newcommand{\clip}{\operatorname{clip}}
\newcommand{\round}{\operatorname{round}}

\title{Digitisation of Rendered Images and Texture Sources}
\author{}
\date{}

\begin{document}
\maketitle

\section{Purpose and common convention}

This note records the image and texture digitisation convention shared by the
render and analysis code.  The durable render product is a floating-point
\texttt{.npy} image.  An 8-bit TIFF is only a visualisation preview.  All
measurement analyses form their requested bit depth directly from the stored
floating-point image using the same function below; higher-bit TIFFs are not a
separate source of truth.

For $B$ bits, define the largest unsigned camera code as
\begin{equation}
  M_B = 2^B-1.
\end{equation}
For a normalised camera intensity $I\in\mathbb{R}$, the common digitiser is
\begin{equation}
  \boxed{
  Q_B(I) = \round\!\left(M_B\,\clip(I,0,1)\right),
  \qquad
  \clip(z,0,1)=\min\{1,\max\{0,z\}\}.}
  \label{eq:camera-digitiser}
\end{equation}
The implementation is \texttt{numpy.rint}: it rounds to the nearest integer,
with exact half-way cases rounded to the even integer.  The output type is
\texttt{uint8} for $B\leq 8$ and \texttt{uint16} otherwise.

Equation~\eqref{eq:camera-digitiser} is implemented centrally as
\texttt{modules.render\_outputs.quantise\_camera}.  It is used for render
previews and for the on-the-fly digitised-error analyses in Experiments 1--3.

\section{Additive speckle images}

Let $c(\bm{x})$ be the raw additive coverage field.  For a camera pixel
$P$, its pixel-box integration---analytic where available, otherwise SSAA---is
\begin{equation}
  \bar c_P = \frac{1}{|P|}\int_P c(\bm{x})\,\mathrm{d}\bm{x}.
  \label{eq:coverage-integral}
\end{equation}
Disk and Gaussian contributions add in $c$; hence overlapping unit-amplitude
blobs can give $\bar c_P>1$.  This raw coverage is intentionally retained in
floating-point additive texture assets and in the float-texture Riley path.
No clipping occurs during texture generation or before pixel integration.

Only after the final pixel value has been integrated is coverage mapped to
normalised image intensity:
\begin{align}
  s_P &= \clip(\bar c_P,0,1),\\
  I_P &= \clip\!\left(I_0+\gamma(1-2s_P),0,1\right).
  \label{eq:coverage-intensity}
\end{align}
Thus zero coverage is the bright level $I_0+\gamma$, and saturated coverage is
the dark level $I_0-\gamma$.  With the current $I_0=0.5$ and $\gamma=0.4$,
these limits are $0.9$ and $0.1$, respectively.  The final camera code is
$Q_B(I_P)$ from Equation~\eqref{eq:camera-digitiser}.

The ordering is important:
\begin{equation}
 \clip\!\left(\frac{1}{|P|}\int_P c(\bm{x})\,\mathrm{d}\bm{x},0,1\right)
 \neq
 \frac{1}{|P|}\int_P\clip(c(\bm{x}),0,1)\,\mathrm{d}\bm{x}
 \quad\text{in general}.
\end{equation}
The code uses the left-hand expression.  It is therefore an additive,
post-integration saturation model rather than the geometric union of opaque
disks.

\section{Eggbox images}

The eggbox is already an intensity function, rather than a coverage field.  In
reference coordinates with pitches $p_x,p_y$, it is
\begin{equation}
 I_{\mathrm{egg}}(x,y)=I_0-\gamma+
 \frac{\gamma}{2}
 \left[1+\cos\!\left(\frac{2\pi x}{p_x}\right)\right]
 \left[1+\cos\!\left(\frac{2\pi y}{p_y}\right)\right].
\end{equation}
The renderer integrates this intensity over the pixel footprint, then applies
$Q_B$ only when an integer camera image is required.  Eggbox float textures
therefore store intensity directly; additive float textures store raw coverage.

\section{Float and uint texture sources}

For an additive float texture, a texel stores the area-averaged raw coverage
$\bar c_T$.  A uint texture of depth $B$ is derived from that same float asset,
not regenerated from the procedural speckle field:
\begin{equation}
 T^{\mathrm{uint}}_B = Q_B\!\left(
 I_0+\gamma\left[1-2\clip(\bar c_T,0,1)\right]\right).
 \label{eq:uint-texture}
\end{equation}
For eggbox textures, whose float value is already intensity,
\begin{equation}
 T^{\mathrm{uint}}_B = Q_B(T^{\mathrm{float}}).
\end{equation}
Consequently, changing the uint source bit depth changes only the quantised
texture source supplied to Riley; it does not repeat analytic texture
integration.  A separate Riley render remains necessary because interpolation
of a quantised source is physically a different operation.

\section{Persisted outputs and analysis}

The standard persisted final-image convention is a normalised floating-point
camera intensity $I_P$ after Equation~\eqref{eq:coverage-intensity} (or the
eggbox equivalent).  The render helpers save this \texttt{float64 .npy} file
and, if absent, derive only an 8-bit TIFF preview using $Q_8$.  For a requested
analysis depth $B$, both candidate and reference float images are independently
converted with $Q_B$ before their integer-code difference is calculated.  This
makes all reported bit-depth metrics reproducible from one canonical floating
image, without depending on a TIFF encoder or pre-existing $B$-bit files.

\end{document}
